\documentclass[12pt]{article}
\usepackage[utf8]{inputenc}
\usepackage[T1]{fontenc}
\usepackage[margin=1in]{geometry}
\usepackage{setspace}
\usepackage{xcolor}
\usepackage[colorlinks=true, urlcolor=blue!60!black, linkcolor=blue!60!black]{hyperref}
\setlength{\parindent}{2em}
\setlength{\parskip}{1.25em}
\renewcommand{\baselinestretch}{1.0}
\begin{document}
\begin{center}
{\LARGE \textit{Palko v.\ Connecticut}, 302 U.S.\ 319 (1937)\par}
\vspace{0.8em}
{\large Case Brief --- Four-Page Edition\par}
\end{center}
\vspace{1em}

\section{Memory Jogger}

Rejecting total incorporation, Justice Cardozo announced the ``ordered liberty'' standard for selective incorporation and held 8--1 that due process did not bar Connecticut from appealing legal errors to retry a murder defendant and seek death, a specific holding overruled in \textit{Benton v.\ Maryland} (1969) (\href{https://supreme.justia.com/cases/federal/us/302/319/}{Justia}).

\section{Detailed Case Facts}

In 1935 Frank Palka (misspelled ``Palko'' in court records) broke into a Bridgeport music store, stole a radio, and fatally shot two police officers while fleeing. At his first trial for first-degree murder, the judge excluded his confession, restricted the prosecution's impeachment cross-examination, and instructed on the difference between first- and second-degree murder. The jury convicted him of second-degree murder, and he received life imprisonment. Under an 1886 statute allowing the State, with the trial judge's permission, to appeal rulings of law in criminal cases, prosecutors appealed. The Connecticut Supreme Court of Errors found prejudicial legal error and ordered a new trial. Over Palko's double-jeopardy objection, the second jury convicted him of first-degree murder, and he was sentenced to death.

\section{Procedural History}

The Supreme Court of Errors reversed the first judgment in 1936 and, after retrial, affirmed the conviction and death sentence in 1937, following its 1894 precedent upholding the state-appeal statute. Palko appealed as of right; the Supreme Court affirmed on December 6, 1937. Palka was executed on April 12, 1938.

\section{Judicial Votes}

Affirmed, \textbf{8--1}.
\begin{itemize}
\item \textbf{Majority:} Justice Cardozo, joined by Chief Justice Hughes and Justices McReynolds, Brandeis, Sutherland, Stone, Roberts, and Black (who would later champion total incorporation in \textit{Adamson}).
\item \textbf{Dissent:} Justice Butler, without opinion.
\end{itemize}

\section{Holding}

The Fourteenth Amendment does not automatically incorporate the first eight amendments against the States. A guarantee is absorbed into due process only if it is ``of the very essence of a scheme of ordered liberty,'' so that ``neither liberty nor justice would exist if [it] were sacrificed.'' A state retrial to correct substantial legal error is not ``a hardship so acute and shocking that our polity will not endure it.'' The Court reserved what due process would require after ``a trial free from error.'' Cardozo's remarks excluding jury trial and self-incrimination were later repudiated (\textit{Duncan}, \textit{Malloy}), but the ``ordered liberty'' framework and his praise of free speech endured.

\section{Analysis of Opinions}

\textbf{Majority (Cardozo).} Assuming such a retrial would violate the Fifth Amendment in a federal case (\textit{Kepner}), Cardozo warned against ``the tyranny of labels'' and rejected the thesis that whatever the Fifth Amendment forbids, the Fourteenth does too. Precedent ran both ways: States could dispense with grand juries (\textit{Hurtado}), juries (\textit{Maxwell v.\ Dow}, \textit{Walker v.\ Sauvinet}), and the self-incrimination privilege (\textit{Twining}), yet the Fourteenth Amendment absorbed the First Amendment freedoms (\textit{Near}, \textit{De Jonge}) and counsel in a capital case (\textit{Powell v.\ Alabama}). The ``rationalizing principle'' is whether a right is implicit in ordered liberty. Juries and the privilege, though valuable, are not indispensable to a ``fair and enlightened system of justice,'' as Continental systems show, while freedom of thought and speech lies on a ``different plane of social and moral values.'' Connecticut was not wearing Palko down after an error-free trial; it sought only ``a trial free from the corrosion of substantial legal error.'' ``The edifice of justice stands,'' Cardozo wrote, ``its symmetry, to many, greater than before.'' The Privileges or Immunities claim was rejected in two sentences.

\textbf{Dissent (Butler).} Butler dissented without opinion, leaving no statement of reasons.

\section{Examples: Future Applications}

\begin{enumerate}
\item \textbf{Same-Side: Charging by Information.} A State charges a felony by prosecutor's information rather than indictment. The Grand Jury Clause is still unincorporated.
\item \textbf{Same-Side: No Civil Juries.} A State replaces civil jury trials with bench trials; under \textit{Palko}, citing \textit{Walker v.\ Sauvinet}, the Seventh Amendment is not essential to ordered liberty.
\item \textbf{Opposite-Side: Prior Restraint.} A State lets officials enjoin newspapers from criticizing judges. Free speech, the ``matrix'' of other freedoms, binds the State.
\item \textbf{Opposite-Side: Retrial After Error-Free Acquittal.} Retrying a defendant because the State disagrees with an acquittal could be the shocking hardship Cardozo reserved, and after \textit{Benton} it is barred outright.
\item \textbf{Fence-Sitter: Shifting Safeguards.} Non-unanimous juries or excessive fines survive \textit{Palko}'s universal test, but under \textit{Duncan} and \textit{McDonald}'s American-tradition test the Court incorporated jury trial, excessive fines (\textit{Timbs}, 2019), and unanimity (\textit{Ramos}, 2020).
\end{enumerate}

\section{Critique}

\textbf{From both flanks.} Total-incorporation originalists (Justice Black's \textit{Adamson} dissent, Akhil Reed Amar, Justice Thomas in \textit{McDonald}) fault \textit{Palko} for using due process rather than the Privileges or Immunities Clause and for letting judges pick rights. Strict federalists see subjective policymaking untethered from text.

\textbf{Civil liberties.} Louis Henkin and Warren Court opinions criticized the abstract ``civilized system'' benchmark as stingy, since Anglo-American adversarial justice relies on jury and double-jeopardy protections; \textit{Benton} rejected \textit{Palko}'s ``watered-down'' double jeopardy. \textbf{Competing values.} Cardozo sought to keep States as laboratories of criminal procedure without sacrificing core freedoms; Palko's execution after a state-initiated retrial shows the cost of that deference.

\section{Key Quotations}

\begin{quote}
``In these and other situations, immunities that are valid as against the federal government by force of the specific pledges of particular amendments have been found to be implicit in the concept of ordered liberty, and thus, through the Fourteenth Amendment, become valid as against the states.''
\end{quote}
\noindent --- Justice Cardozo, for the Court, 302 U.S.\ at 324--25.

\begin{quote}
``If the Fourteenth Amendment has absorbed them, the process of absorption has had its source in the belief that neither liberty nor justice would exist if they were sacrificed. This is true, for illustration, of freedom of thought, and speech. Of that freedom one may say that it is the matrix, the indispensable condition, of nearly every other form of freedom.''
\end{quote}
\noindent --- Justice Cardozo, for the Court, 302 U.S.\ at 326--27.

\begin{quote}
``The right to trial by jury and the immunity from prosecution except as the result of an indictment may have value and importance. Even so, they are not of the very essence of a scheme of ordered liberty.''
\end{quote}
\noindent --- Justice Cardozo, for the Court, 302 U.S.\ at 325.

\end{document}
